\documentclass[protokoll]{dvd}

\KOMAoptions{
    headwidth = 18cm,
    footwidth = 18cm,
}

% Styrelsen
\def\ordf{Daniell Cole}
\def\vice{\vak}
\def\kassa{Owais Tabbaa}
\def\samo{Angelika Hagberg}
\def\sekr{\vak}
\def\vak{\emph{vakant}}

% mötesinfo

\def\nr{4}
\def\yr{2026}
\def\ymd{\yr-10-15}
\def\starttime{11:11}


\begin{document}

\title{Styrelsemöte \nr}
\subtitle{\yr}
\author{Styrelsen}
\date{\ymd}


\textbf{Datum:} \csname @date\endcsname\\
\textbf{Tid:} \starttime\\
\textbf{Plats:} Styrelserummet\\
\textbf{Styrelsemedlemmar:}
\begin{närvarande_förtroendevalda}
    \förtroendevald{Ordförande}{\ordf}{Ja}
    \förtroendevald{Vice Ordförande}{\vice}{---}
    \förtroendevald{Kassör}{\kassa}{Ja}
    \förtroendevald{SAMO}{\samo}{Ja}
    \förtroendevald{Sekreterare}{\sekr}{---}
\end{närvarande_förtroendevalda}

%\textbf{Övriga medlemmar:}
%\begin{närvarande_medlemmar}
%\medlem{Ingen}[]
%\end{närvarande_medlemmar}
\section{Öppnande av möte}

Mötet planeras öppnas av \ordf\ kl \starttime

Mötet öppnades 11:30

%Mötet öppnades 
\section{Runda bordet}

Runda bordet innebär att varje person berättar hur de känner sig.
Man kan till exempel berätta att man är stressad på grund av en inlämning, irriterad på sin granne, eller bara väldigt glad därför att man ligger i fas med plugget.

\section{Formalia}

\subsection{Styrelsens beslutbarhet}

\blockquote[7 kap. 5 \S~första stycket i stadgan][]{%
    Styrelsen är endast beslutsmässig då samtliga styrelsemedlemmar har fått kallelsen till styrelsemötet och minst hälften av styrelsemedlemmarna är närvarande. Ordförande eller vice ordförande måste vara närvarande när beslut tas.
}

\blockquote[7 kap. 11 \S~första stycket i stadgan]{
Styrelsens sammanträden måste kallas minst två läsdagar innan sammanträdet. }

\blockquote[7 kap. 11 \S~sista stycket i stadgan]{
Kallelse till sammanträde får endast ske när protokollet från tidigare sammanträde är justerat och tillgängligt för medlemmarna.}

\subsubsection*{Förslag}

\begin{attsatser}
    \item Mötet kallades 2026-10-06 samt gjordes protokollet för styrelsemöte 3 tillgängligt genom verktyget på hemsidan på samma dag.
    Vilket innebär att styrelsen har uppnått kraven i 7 kap. 5 § första stycket i stadgan och är därmed beslutbar.
\end{attsatser}
\subsubsection*{Beslut}
\begin{attsatser}
    \item förslaget till beslut bifalles
\end{attsatser}


\subsection{Fastställande av mötesschema}

För att styrelsen ska kunna fatta ett styrelsebeslut eller protokollföra en diskussion behöver punkten i mötesschemat där styrelsen ska fatta beslut vara inlagd eller föras in i mötesschemat senast vid den här punkten.

\subsubsection*{Förslag}

\begin{attsatser}
    \item 
    mötesschemat fastställs utan några förändringar.
    %punkt läggs till på mötesschema
\end{attsatser}
\subsubsection*{Beslut}
\begin{attsatser}
    \item förslaget till beslut bifalles
\end{attsatser}

\subsection{Val av mötessekreterare}
Vi behöver någon som antecknar. 

\subsubsection*{Förslag}
\begin{attsatser}
    \item Angelika väljs till mötessekreterare
\end{attsatser}
\subsubsection*{Beslut}
\begin{attsatser}
    \item förslaget till beslut bifalles
\end{attsatser}

\subsection{Val av protokolljusterare}

Protokolljusterare har till uppgift att kontrollera att protokollet i slutändan reflekterar de faktiska besluten och diskussionerna som fördes under mötet.
Utöver protokolljusteraren så ska mötesordförande och mötessekreteraren signera protokollet.
Vid styrelsemöten ska det endast vara en justerare.
Mötesordförande och mötessekreteraren kan inte vara justerare.

\subsubsection*{Förslag}
\begin{attsatser}
    \item Owais är justerare
\end{attsatser}
\subsubsection*{Beslut}
\begin{attsatser}
    \item förslaget till beslut bifalles
\end{attsatser}

\section{Rapport} 
\subsection{Ordförande}
Aspinfomötet hade ett bra antal personer som dök upp, utöver det har jag inte gjort så mycket.


Nästa SUS-utbildningstillfälle har annonserats och vi har därmed möjlighet att skicka två pers till det. 

Vi har nu officiellt fått godkänd kårföreningsstatus. 

NaTsek-sektionsstyrelsemöte nr 1 kommer ske idag (samma dag som detta möte), jag hade först tänkt gå på det men jag har föreläsning tills 17 och behöver vara på flygplatsen kl 08:30 imorgon (fredag 16:e). 

\subsection{Kassör}
Kollat igenom excelarket, det mesta är utbetalt, väntar på formulär för sista betalningar. Vi ligger bra till och har ca 28 000 tusen kvar på kontot när allt är klart. DVRK
skickade äskan för strutbygget för sent (efter att strutbygget ägt rum) och kunde därför inte få
den godkänt. Ska prata med Matti angående äskningar och dess rapport efter tentaveckan.

Äskningar under 200 för fika godkändes åt concats och dvrk.

\subsection{SAMO}
Inget att rapportera

\newpage


\section{Diskussionspunkter}
\subsection*{Förslag på stadgaändringar}

\subsubsection*{Förslag 1}
Minimum 3 medlemmar i styrelsen: 
\begin{attsatser}
    \item i början av 7 kap. 3 § i stadgan införa 
    \begin{quote}
        Styrelsen kan som minst bestå av tre medlemmar varav en skall vara divisionsordförande och en divisionskassör. 
    \end{quote}
\end{attsatser}

\subsubsection*{Förslag 2}
Krav från Göta, styrelsemedlem behöver vara ordinarie medlem i Göta studentkår:
\begin{attsatser}
    \item ändra första stycket i 7 kap. 2 § från
    \begin{quote}
        Styrelsens medlemmar väljs av divisionsstämman. Mandatperioden är densamma som räkenskapsåret. Valbar till styrelsen är medlem i divisionen.
    \end{quote}
    till
    \begin{quote}
        Styrelsens medlemmar väljs av divisionsstämman. Mandatperioden är densamma som räkenskapsåret. Valbar till styrelsen är medlem i divisionen med ordinarie medlemskap i Göta studentkår.
    \end{quote}
\end{attsatser}

\subsubsection*{Förslag 3}
Ett till krav från Göta, sammanfatta vilka rättigheter föreningsmedlemmar har:
\begin{attsatser}
    \item införa i 4 kap, 4 § av stadgan 
    \begin{quote}
        Medlemskap medför: 
        \begin{itemize}
            \item rätten att kalla till divisionsstämma (6 kap, 1 §),
            \item närvaro-, yttrande, förslags- och rösträtt på divisionsstämman (6 kap, 5 §),
            \item rätt att bli medlem i styrelsen (7 kap, 2 §),
            \item närvarorätt på styrelsemöten (7 kap, 6 §),
            \item samt rätten att bli medlem i kommittéer (8 kap, 4 §). 
        \end{itemize}
    \end{quote}
\end{attsatser}

\subsubsection*{Förslag 4}
Den stora jag främst vill ha feedback på, ny roll i styrelsen; Styrelseledamot:
\begin{attsatser}
    \item ändra 7 kap, 1§ från:
    \begin{quote}
    Styrelsen består av
		\begin{itemize}
			\item ordförande (även kallad divisionsordförande);
			\item vice ordförande;
			\item sekreterare;
			\item kassör (även kallad divisionsskassör);
			\item SAMO; och
            \end{itemize}
    \end{quote}
    till:
    \begin{quote}
        Styrelsens ordinarie medlemmar består av 
		\begin{itemize}
			\item ordförande (även kallad divisionsordförande);
			\item vice ordförande;
			\item sekreterare;
			\item kassör (även kallad divisionsskassör);
			\item SAMO; och
            \end{itemize}
            ytterligare styrelsemedlemmar består av
            \begin{itemize}
                \item ledamöter (även kallat    styrelseledamöter).
            \end{itemize}
    \end{quote}
    
    \item införa i 7 kap. 2 §:
    \begin{quote}
        Ordinarie styrelsemedlem kan inte samtidigt vara styrelseledamot.
    \end{quote}
    
    \item uppdatera 7 kap. 3 § från:
    \begin{quote}
        Styrelsen kan högst bestå av fem medlemmar.
    \end{quote}
    till: 
    \begin{quote}
        Styrelsen kan högst bestå av fem ordinarie medlemmar. 
    \end{quote}
    
    \item 
    ändra första stycket i 7 kap. 5 § från:
    \begin{quote}
        Styrelsen är endast beslutsmässig då samtliga styrelsemedlemmar har fått kallelsen till styrelsemötet och minst hälften av styrelsemedlemmarna är närvarande.
        Ordförande eller vice ordförande måste vara närvarande när beslut tas.
    \end{quote}
    till:
    \begin{quote}
        Styrelsen är endast beslutsmässig då samtliga styrelsemedlemmar har fått kallelsen till styrelsemötet och minst hälften av styrelsens ordinarie medlemmar är närvarande.
        Ordförande eller vice ordförande måste vara närvarande när beslut tas.
    \end{quote}
\end{attsatser}

\begin{quote}
    \textit{Vi håller med om punkterna, ska ta hjälp av tidigare styrelsemedlemmar för att se över formuleringen och ta emot deras åsikter.}
\end{quote}

\subsection*{Nästa stämma}
Bara en stämma i december?
\\
\begin{quote}
    \textit{Nästa styrelsemöte ska vi bestämma datumet på nästa stämma, bara en stämma innan årsskiftet, blir i december.}
\end{quote}


\subsection*{Luciafika/Jubord}
\begin{quote}
    \textit{Ska använda pengar som är över från insparken. Har pengar till ett julbord om vi får in volontärer som kan hjälpa till. Diskutera vidare detta i möten framöver.}
\end{quote}

\section{Besultspunkter}

\subsection{SUS info}
\begin{quote}
    Dag 1 onsdag 11:e november 11:00-16:00 (20 min senare slut för chalmerister) i Scaniasalen på Chalmers
    Dag 2 fredag 13:e november 10:00-12:00 i Scaniasalen på Chalmers
    ...
    max två per förening, och priset är 300 kr.
\end{quote}

Med att jag kommer vara borta i en vecka skulle det vara bra om  någon annan tar och skickar ut denna info samt hanterar anmälan. 

Förra året skickade jag ut infon och bad de intresserade att kontakta mig och sedan skickade jag länken till anmälningsformuläret privat till dem. 

\subsubsection*{Förslag till beslut}
\begin{attsatser}
    \item  Daniell skickar ut infon om två veckor. 
\end{attsatser}
\subsubsection*{Beslut}
\begin{attsatser}
    \item Attsats 1 bifalles.
\end{attsatser}

\subsection{Äskan: Fika - Concats pluggstuga}
\begin{quote}
Hej!
Vi i ConCats tänker hålla en pluggstuga den 21:a och vi skulle förvänta ganska många då så vill vi äska för 500 kr om det kommer 10+ personer.

Mvh Owais ``Mirage''\\
ConCats
\end{quote}

10+ är lite vagt, förmodligen 15-20. 500 kr är en rimlig mängd till det antalet och inför de nya ettornas första riktiga tenta. 

\subsubsection*{Förslag till beslut}
\begin{attsatser}
    \item äskan på 500 kr åt fika till pluggstuga 21-10-2026 godkänns. 
\end{attsatser}
\subsubsection*{Beslut} 
\begin{attsatser}     
\item Attsats 1 bifalles. 
\end{attsatser}

\subsection{Äskan: Idrottshall - DVRK 5-kamp}
\begin{quote}
Hej vi i dvrk vill hålla i 5 kamp den 13/11. Vi vill hålla eventet i Landala idrottshall och det kostar 1282kr att hyra den för 2h. Det förväntas komma runt 15 personer på det eventet.

Vänligen,\\
Halimo ``Lost\&\&Found''
\end{quote}

Förslag att vi tar från pengarna som kommit in från kårkällar'n istället för CSE pengar.

\subsubsection*{Förslag till beslut}
\begin{attsatser}     
\item äskan på 1282 kr för idrottshall till 5-kamp 13-11-2026 godkänns.
\end{attsatser}
\subsubsection*{Beslut} 
\begin{attsatser}     
\item Attsats 1 bifalles. 
\end{attsatser}

\section{Per Capsulam-beslut}
\blockquote[]{
När per capsulam-beslut har fattats måste beslutet dokumenteras,\\
\ \cdot\ \ på divisionens Discordserver som är tillgänglig alla studenter;\\
\ \cdot\ \ i protokollet för det första styrelsesammanträdet efter beslut har fattats.\\

Dokumentationen ska innehålla\\
\ \cdot\ \ vilka att-satser som har bifallits;\\
\ \cdot\ \ omröstningsperioden.\\

Om dessa krav inte uppfylls upphävs beslutet.
}

\textit{Inga per capsulam-beslut har tagits.}

% \subsubsection*{Bifallna attsatser}
% \subsubsection*{Omröstningsperioden}

\newpage
\section{Avslutande av möte}

\subsection{Nästa möte}
2026-11-02 - 12:00

\subsection{Mötets avslutande}
Mötet avslutades kl 12:35 

\styrelsesignaturer

\end{document}